% =====================================================================
%  Informe de auditoria de la tesis EduFEM
%  Generado el 2026-06-03. Documento independiente (no forma parte de la
%  tesis). Compila con pdflatex (MiKTeX).
% =====================================================================
\documentclass[11pt,a4paper]{article}

\usepackage[utf8]{inputenc}
\usepackage[T1]{fontenc}
\usepackage[spanish,es-tabla]{babel}
\usepackage{lmodern}
\usepackage[a4paper,margin=2.3cm]{geometry}
\usepackage{amsmath,amssymb}
\usepackage{bm}
\usepackage{booktabs}
\usepackage{array}
\usepackage{longtable}
\usepackage{enumitem}
\usepackage{xcolor}
\usepackage[most]{tcolorbox}
\usepackage{hyperref}
\usepackage{microtype}
\usepackage{parskip}
\usepackage{fancyhdr}

\hypersetup{
  colorlinks=true,
  linkcolor=blue!60!black,
  urlcolor=blue!60!black,
  pdftitle={Informe de auditoria - Tesis EduFEM},
  pdfauthor={Auditoria asistida},
}

\definecolor{cAlta}{HTML}{C0392B}
\definecolor{cMedia}{HTML}{D68910}
\definecolor{cBaja}{HTML}{707B7C}
\definecolor{cOk}{HTML}{1E8449}
\definecolor{cAzul}{HTML}{0D6EFD}

\newcommand{\sevalta}{\textcolor{cAlta}{\textbf{[ALTA]}}}
\newcommand{\sevmedia}{\textcolor{cMedia}{\textbf{[MEDIA]}}}
\newcommand{\sevbaja}{\textcolor{cBaja}{\textbf{[BAJA]}}}
\newcommand{\sevok}{\textcolor{cOk}{\textbf{[OK]}}}
\newcommand{\fileref}[1]{\texttt{\small #1}}

\tcbuselibrary{skins,breakable}
\newtcolorbox{execbox}{
  colback=blue!4, colframe=cAzul, boxrule=0.8pt, arc=3pt,
  left=10pt,right=10pt,top=8pt,bottom=8pt, breakable,
  title=\textbf{Resumen ejecutivo}, fonttitle=\color{white}\bfseries,
  colbacktitle=cAzul
}
\newtcolorbox{notabox}{
  colback=gray!5, colframe=gray!50!black, boxrule=0.6pt, arc=2pt,
  left=8pt,right=8pt,top=6pt,bottom=6pt, breakable
}

\pagestyle{fancy}
\fancyhf{}
\fancyhead[L]{\small Auditoría — Tesis EduFEM}
\fancyhead[R]{\small\thepage}
\renewcommand{\headrulewidth}{0.4pt}

\setlist[itemize]{leftmargin=1.4em, itemsep=2pt, topsep=3pt}

\title{\vspace{-1.2cm}\textbf{Informe de auditoría}\\[2pt]
\large Tesis: \textit{Desarrollo de software educativo de elementos finitos para el análisis estructural empleando el lenguaje de programación Python} (EduFEM)}
\author{Revisión técnica y editorial asistida}
\date{3 de junio de 2026}

\begin{document}
\maketitle
\thispagestyle{fancy}

\begin{execbox}
La tesis se encuentra en un estado \textbf{sólido y apto para defensa}. La arquitectura
argumental es madura y cierra el lazo problema\,--\,objetivo\,--\,hipótesis\,--\,verificación\,--\,conclusiones,
con correspondencia uno a uno entre objetivos específicos y conclusiones. El rigor técnico y la
consistencia numérica son altos: los datos cuadran entre capítulos y la memoria de cálculo del
anexo es trazable, cifra por cifra, con la salida del motor. El marco metodológico está bien
preparado para un tribunal que aplique los enfoques de Álvarez de Zayas (problema, objeto, campo,
objetivos, hipótesis) y de Hernández-Sampieri (enfoque, alcance, diseño).

\medskip
La auditoría en siete dimensiones reunió \textbf{43 hallazgos}: \sevalta~2, \sevmedia~15 y
\sevbaja~26. \textbf{Ningún hallazgo compromete la validez de los resultados}; los dos de severidad
alta son de estilo (fatiga léxica por repetición), no de contenido. La intervención previa —cuatro
mejoras de contenido y una revisión adversarial (\autoref{sec:intervencion})— ya está incorporada y
el documento compila limpio (122~páginas, sin referencias ni citas sin resolver). Las observaciones
de la auditoría que recaían sobre el texto recién añadido se corrigieron en el acto; el resto se
entrega en este informe para decisión del autor.

\medskip
\noindent\textbf{Prioridades sugeridas:}
\begin{enumerate}[leftmargin=1.6em,nosep]
  \item Corregir la única inconsistencia aritmética genuina: el error frente a Euler-Bernoulli
  declara \(1{,}82\%\) cuando el cálculo da \(1{,}85\%\) (\fileref{04\_resultados.tex}); verificar
  contra el guion \fileref{vv\_timoshenko}.
  \item Podar la redundancia transversal: limitaciones repetidas en tres capítulos y solapamiento de
  la descripción de módulos entre el capítulo de diseño y la sección de resultados.
  \item Diversificar el léxico sobre-usado (``canal de cálculo'', ``caja negra'', ``paso a paso'').
  \item Reforzar el flanco metodológico: distinguir los dos planos de variables (numérico y
  pedagógico) y dotar a la pregunta PI-2 de un criterio de cumplimiento auditable.
  \item Cerrar las afirmaciones comparativas sin fuente sobre productos de terceros y ampliar
  mínimamente la base de literatura pedagógica.
\end{enumerate}
\end{execbox}

\section{Alcance y método de la auditoría}

Este informe acompaña una intervención de mejora sobre la tesis y presenta una auditoría
integral de su estado resultante. La revisión se organizó en dos fases. En la primera se
aplicaron cuatro cambios de contenido solicitados por el autor y se sometieron a una
\emph{revisión adversarial} por lentes independientes (rigor técnico, registro, higiene
\LaTeX, consistencia interna y citación), cuyas observaciones se incorporaron al texto. En
la segunda se auditó la tesis completa —los trece archivos fuente \texttt{.tex}, la
bibliografía y los recursos de figuras— a lo largo de siete dimensiones, cada una a cargo de
un revisor especializado independiente.

La severidad de cada hallazgo se clasifica como \sevalta\ (debe corregirse antes de la
defensa), \sevmedia\ (recomendable corregir; fortalece la solidez) o \sevbaja\ (pulido
opcional). Se distingue, además, entre defectos genuinos y \emph{decisiones de autor}
legítimas que un evaluador podría cuestionar pero que son defendibles. La tesis compila de
forma limpia (122~páginas, sin referencias ni citas sin resolver) en su estado actual.

\section{Intervención aplicada (junio de 2026)}
\label{sec:intervencion}

La intervención partió de cuatro observaciones del autor. La \autoref{tab:intervencion}
resume cada una, su tratamiento y los archivos afectados.

\subsection*{1. Justificación del dominio bidimensional como puente 1D\,$\leftrightarrow$\,3D}

Se incorporó en la Introducción (sección \emph{Alcance y limitaciones}) un párrafo que
reencuadra la restricción a dos dimensiones como una \emph{decisión pedagógica}, no como una
mera limitación de alcance. El argumento sitúa el continuo plano en el punto intermedio entre
el análisis unidimensional de barras —insuficiente para ilustrar el mapeo isoparamétrico, el
Jacobiano, la matriz $\bm{B}$ y la cuadratura en dos direcciones— y el tridimensional, cuyo
aparato algebraico ($\bm{D}$ de $6\times6$, $\bm{B}$ de $6\times24$ y rigidez de $24\times24$
en un hexaedro de ocho nodos) volvería impracticable el seguimiento paso a paso que constituye
el núcleo didáctico de la herramienta. El plano conserva matrices de tamaño legible
($\bm{D}$ de $3\times3$, $\bm{B}$ de $3\times8$ en el Q4) y, a la vez, una formulación
generalizable a 3D, de modo que opera como puente: lo aprendido en 2D se extiende de forma
natural al caso tridimensional, planteado como línea de continuación.

\subsection*{2. Anclaje teórico del bloqueo por cortante}

El \emph{bloqueo por cortante} (\emph{shear-locking}) se mencionaba en la Introducción, en los
Resultados (membrana de Cook) y en las Conclusiones, pero carecía de un tratamiento teórico
propio en el Marco teórico, y la referencia al bloqueo volumétrico era una promesa
incumplida (``se discute más adelante''). Se añadió una sección dedicada —\emph{Fenómenos
numéricos en elementos de bajo orden: bloqueo y modos espurios}— que define con respaldo
bibliográfico el bloqueo por cortante, el bloqueo volumétrico y los modos espurios
(\emph{hourglass}), y discute las técnicas correctivas (SRI, $\bar{\bm{B}}$) que la herramienta
no implementa. La promesa colgante se reconectó mediante referencia cruzada, y las menciones
dispersas ahora apuntan a esta sección.

\subsection*{3. Corrección del dato sobre ED-Elas2D}

La tesis describía al antecedente ED-Elas2D como software ``en inglés''. La verificación del
material original confirma que es un programa desarrollado en la Universitat Politècnica de
Catalunya, de origen español y con versiones en español; el artículo de revista que lo
documenta está en inglés porque la publicación es internacional, pero el idioma no es un rasgo
diferenciador frente a EduFEM. Se corrigió el dato (incluida la celda de la tabla comparativa,
de ``Inglés'' a ``Español'') y se reenfocó la diferenciación hacia los atributos que sí son
genuinos y defendibles: carácter libre y de código abierto (frente a software académico
cerrado), exposición interactiva del cálculo por elemento, generación automática de una
memoria de cálculo trazable y batería de verificación y validación documentada. Esta
corrección protege la tesis de una objeción factual del tribunal.

\subsection*{4. Realce de la memoria de cálculo paso a paso como aporte distintivo}

Se reforzó —en el Resumen, la Introducción, los Resultados y las Conclusiones— la generación
automática de una memoria de cálculo que reconstruye el procedimiento con sustitución numérica
paso a paso sobre el modelo del propio usuario, presentándola como un \emph{compañero de
verificación}: un registro que el estudiante puede rehacer a mano o en una hoja de cálculo y
contrastar con el programa, elemento por elemento. Para mantener el rigor, las afirmaciones de
rareza se acotaron al relevamiento de antecedentes del propio trabajo (``no se documenta entre
los antecedentes revisados'') en lugar de aseverar una unicidad no medida.

\subsection*{Revisión adversarial y consistencia}

Las ediciones se sometieron a revisión por lentes independientes. El hallazgo de mayor impacto
fue una \textbf{inconsistencia crítica derivada del cambio 3}: el ``problema general'' del
diseño metodológico seguía afirmando la ``ausencia de un entorno educativo en español'', lo que
el reconocimiento de ED-Elas2D como software español volvía falso. Se reancló el problema en la
combinación inédita (libre y de código abierto, en español, con memoria de cálculo trazable),
coherente ya con el marco teórico. Otras correcciones incorporadas: citas localizadas en cada
bloqueo, respaldo bibliográfico del párrafo del dominio 2D, atenuación de afirmaciones
comparativas, reducción de redundancia en la sección de \emph{hourglass} y una aclaración de que
la tensión de von Mises se calcula en su forma plana (consistente con el motor de cálculo).

\begin{table}[h]
\centering
\caption{Síntesis de la intervención aplicada.}
\label{tab:intervencion}
\small
\renewcommand{\arraystretch}{1.25}
\begin{tabular}{@{}p{4.0cm}p{6.2cm}p{4.6cm}@{}}
\toprule
\textbf{Observación} & \textbf{Tratamiento} & \textbf{Archivos} \\
\midrule
Dominio 2D como puente 1D--3D & Párrafo de justificación pedagógica con dimensiones de matrices y citas & \fileref{01\_introduccion.tex} \\
Bloqueo por cortante sin anclaje teórico & Nueva sección de fenómenos numéricos (bloqueo y \emph{hourglass}) + referencias cruzadas & \fileref{02\_marco\_teorico.tex}, \fileref{05\_conclusiones.tex} \\
ED-Elas2D no es ``en inglés'' & Corrección del dato + reenfoque de la diferenciación & \fileref{02\_marco\_teorico.tex} \\
Memoria de cálculo como aporte único & Realce en cuatro secciones, con afirmaciones acotadas & \fileref{00\_resumen.tex}, \fileref{01\_introduccion.tex}, \fileref{04\_resultados.tex}, \fileref{05\_conclusiones.tex} \\
Consistencia (revisión) & Reanclaje del problema general; citas, von Mises, redundancia & \fileref{02b\_diseno\_metodologico.tex}, \fileref{02\_marco\_teorico.tex} \\
\bottomrule
\end{tabular}
\end{table}

\section{Auditoría por dimensiones}

Cada dimensión fue auditada por un revisor independiente sobre el estado actual de la tesis. Se
reportan las fortalezas verificadas, los hallazgos con su severidad y ubicación, y un veredicto.
Cuando una observación recaía sobre el texto añadido en esta intervención y se corrigió de
inmediato, se marca \textit{(resuelto)}.

\subsection{Estructura, argumentación y coherencia global}

\textbf{Fortalezas.} La cadena problema\,$\to$\,objetivo\,$\to$\,hipótesis\,$\to$\,resultados está bien
trabada: la ``doble brecha'' (abstracción de los textos / opacidad del software) se enuncia al inicio
y se recupera explícitamente en las conclusiones (``el trabajo regresa así a la contradicción que lo
originó''). La correspondencia objetivos\,$\leftrightarrow$\,conclusiones es uno a uno y verificable, sellada
por la matriz de consistencia. El manejo de la hipótesis ``por diseño'' es honesto y consistente de
extremo a extremo. El resumen refleja fielmente el contenido y sus cifras coinciden con las tablas.

\textbf{Hallazgos.}
\begin{itemize}
  \item \sevmedia\ \fileref{01\_introduccion.tex}, \fileref{04\_resultados.tex} y \fileref{05\_conclusiones.tex} — La misma lista de limitaciones (no 3D, no no-lineal, solo Q4/Q9, bloqueo no mitigado, sin validación con estudiantes) aparece tres veces con redacción muy próxima. \textit{Recomendación:} en Resultados conservar solo las limitaciones que emergen de la evidencia medida (escalabilidad de \fileref{tab:tiempos}, bloqueo de \fileref{tab:cook}) y remitir el resto a la Introducción.
  \item \sevmedia\ \fileref{04\_resultados.tex} (``Resultados del software y cobertura'') solapa con la descripción de módulos del capítulo de diseño (\fileref{tab:modulos}): describe el producto, no un resultado medido. \textit{Recomendación:} reencuadrar como evidencia de cumplimiento del objetivo de cobertura del canal de cálculo y remitir a la tabla, sin re-describir cada módulo.
  \item \sevmedia\ Transversal — La metáfora ``caja negra'' (7 apariciones) y la ``doble brecha'' se reiteran en resumen, introducción, marco y conclusiones. \textit{Recomendación:} fijar la metáfora en el planteamiento y el cierre, y sustituirla por variantes en el cuerpo intermedio.
  \item \sevbaja\ \fileref{tab:resumen-q4q9} — La fila ``Timoshenko $\sigma_x$'' deja Q4 con ``---'', lo que puede leerse como dato faltante (la validación se hizo solo con Q9). \textit{Recomendación:} sustituir por ``(no evaluado; malla Q9)''.
  \item \sevbaja\ \fileref{01\_introduccion.tex} (Estructura del documento) — Afirma ``tres capítulos'' mientras el lector percibe cuatro bloques (marco, metodológico, diseño, resultados). \textit{Recomendación:} una frase que aclare que el diseño metodológico abre el capítulo de diseño e implementación.
  \item \sevbaja\ Completitud — No hay un caso de aplicación de ingeniería civil concreto (una ménsula, un muro) interpretado como lo haría un ingeniero; el salto del banco de prueba abstracto al aula queda implícito. \textit{Recomendación (opcional):} un párrafo breve con la lectura ingenieril de un contorno de tensión reconocible.
  \item \sevbaja\ \fileref{04\_resultados.tex} y \fileref{05\_conclusiones.tex} — La validación de la hipótesis se repite casi literal entre ambos. \textit{Recomendación:} sintetizar en conclusiones y remitir a la sección de validación.
\end{itemize}

\textbf{Veredicto.} Arquitectura argumental sólida y madura; la prioridad es podar la redundancia
transversal. Ningún hallazgo es de severidad alta; la estructura está lista para defensa con ajustes
editoriales menores.

\subsection{Rigor técnico FEM y consistencia numérica}

\textbf{Fortalezas.} La coherencia numérica es notable y verificable: cada cifra de la memoria de
cálculo (Anexo F) reproduce la salida del motor (matriz $\bm{D}$, Jacobiano y su determinante, matriz
de extrapolación de Barlow, tensiones principales y de von Mises, equilibrio de reacciones
$\sum R_y = 1000 = P$). Las tasas de convergencia del MMS se reproducen al recalcular el logaritmo del
cociente de errores (Q4: $\mathcal{O}(h^{2})$ en $L^2$, $\mathcal{O}(h^{1})$ en $H^1$; Q9:
$\mathcal{O}(h^{3})$ y $\mathcal{O}(h^{2})$). La formulación teórica es correcta y sigue el orden del
canal de cálculo real. La honestidad metodológica (nota del valor redondeado de Cook, advertencia de
que la cáscara de SAP2000 no es el continuo de tensión plana) es valiosa.

\textbf{Hallazgos.}
\begin{itemize}
  \item \sevmedia\ \fileref{04\_resultados.tex} — Inconsistencia aritmética: el texto afirma que con la referencia de Euler-Bernoulli ($1{,}9976$~cm) ``el error aparente ascendería al $1{,}82\%$'', pero $(2{,}03460-1{,}9976)/1{,}9976 = 1{,}85\%$. \textit{Recomendación:} corregir a $1{,}85\%$ (o ajustar la referencia a $1{,}9982$~cm); verificar cuál produce el guion \fileref{vv\_timoshenko}. \textbf{Es el único error numérico genuino del documento.}
  \item \sevbaja\ \fileref{04\_resultados.tex} — $217\,370$~kgf/cm$^2 \approx 21{,}32$~GPa, no $21{,}33$. \textit{Recomendación:} corregir a $\approx 21{,}32$~GPa.
  \item \sevbaja\ Deflexión $0{,}2633\%$ frente a $0{,}2632\%$ recalculado de los valores redondeados (el motor usa precisión completa; admisible). \textit{Recomendación:} ninguna acción obligatoria.
  \item \sevbaja\ \fileref{02\_marco\_teorico.tex} — La forma plana de von Mises ($\sigma_3=0$) puede subestimar la tensión equivalente en deformación plana. \textit{(resuelto: se añadió esa aclaración en la nota de la ecuación).}
\end{itemize}

\textbf{Veredicto.} Rigor técnico y consistencia numérica sólidos; los datos cuadran entre capítulos y
la memoria de cálculo es trazable cifra por cifra con el motor. La prioridad es corregir el
$1{,}82\%$ frente a $1{,}85\%$; las dos cifras de redondeo son cosméticas.

\subsection{Marco metodológico y alineación con el tribunal}

\textbf{Fortalezas.} El diseño teórico clásico está explícito y diferenciado (problema científico como
interrogante, objeto de estudio y campo de acción con su relación de inclusión enunciada
literalmente). La hipótesis de diseño está blindada en tres lugares concéntricos y consistentes, y se
traduce a la nomenclatura clásica de variable independiente/dependiente, gesto que anticipa la
objeción del tribunal. La matriz de consistencia es trazable de punta a punta y honesta sobre su
propio hueco (el quinto objetivo, instrumental). Los sistemas de control, repetición y validación se
fundan correctamente en el determinismo del motor.

\textbf{Hallazgos.}
\begin{itemize}
  \item \sevmedia\ \fileref{02b\_diseno\_metodologico.tex} (\fileref{tab:variables} vs. hipótesis) — Conviven dos juegos de variables: las del experimento numérico (tipo de elemento, exactitud) y las de la hipótesis de diseño (transparencia, apoyo a la comprensión), sin un puente que aclare que operan en planos distintos. \textit{Recomendación:} una frase bisagra que explicite los dos niveles (verificación numérica vs.\ diseño pedagógico).
  \item \sevmedia\ \fileref{02b\_diseno\_metodologico.tex} (Población y casos de estudio) — Se llama ``población/muestra'' al universo de problemas-test, no a sujetos; un tribunal formado en Hernández-Sampieri puede objetar la disonancia con el objeto declarado (proceso de enseñanza-aprendizaje). \textit{Recomendación:} una oración que justifique el desplazamiento conceptual y lo conecte con la validación pedagógica futura.
  \item \sevmedia\ PI-2 se verifica ``por diseño'' sin criterio de aceptación explícito, a diferencia de PI-1 y PI-3 (con umbrales numéricos). \textit{Recomendación:} declarar un criterio de diseño auditable —una checklist etapa del canal\,$\leftrightarrow$\,módulo que la expone, con base en \fileref{tab:modulos}—.
  \item \sevbaja\ La caracterización metodológica (propositivo/aplicada-tecnológica/cuantitativo) está duplicada casi literal en introducción y capítulo metodológico. \textit{Recomendación:} sintetizar en la introducción y remitir.
  \item \sevbaja\ \fileref{tab:consistencia} — La columna ``Variable/indicador'' usa para algunos objetivos atributos de diseño no medibles. \textit{Recomendación:} renombrar a ``Variable/indicador o atributo de diseño''.
  \item \sevbaja\ Estructura ``tres capítulos'' con el diseño metodológico dentro del capítulo de diseño; un tribunal ortodoxo espera un capítulo metodológico propio. \textit{Recomendación:} no cambiar la estructura; añadir una frase que justifique la integración como unidad propia de un trabajo tecnológico.
\end{itemize}

\textbf{Veredicto.} Marco metodológico sólido y bien blindado para un tribunal UATF. La prioridad es
reconciliar los dos juegos de variables (numéricas vs.\ pedagógicas) y dotar a PI-2 de un criterio de
cumplimiento auditable.

\subsection{Redacción, estilo académico y naturalidad}

\textbf{Fortalezas.} Voz impersonal y registro académico consistentes; sin coloquialismos ni primera
persona. Terminología MEF canónica impecable (GDL, MEF, restricción, tensión/deformación/desplazamiento;
las únicas apariciones de ``FEM''/``DOF'' están dentro de código). Prosa teórica clara y bien
encadenada. Ausencia de las muletillas de IA más delatoras.

\textbf{Hallazgos.}
\begin{itemize}
  \item \sevalta\ Transversal — Sobreuso de ``canal de cálculo'' (unas 28 apariciones con sus variantes ``cadena de cálculo/del método''). Repetido con esa densidad, suena a fórmula. \textit{Recomendación:} alternar (``el procedimiento del MEF'', ``la secuencia de cálculo'', ``el flujo de operaciones'') y reducir a unos 12--15 usos.
  \item \sevalta\ Transversal — El binomio ``caja negra que oculta\,/\,paso a paso que el estudiante puede replicar'' se repite casi verbatim en resumen, introducción, resultados y conclusiones (``caja negra'' 7 veces, ``paso a paso'' unas 17). \textit{Recomendación:} variar ``paso a paso'' (``con sustitución numérica detallada'', ``operación por operación'') y reformular el binomio en al menos dos de los cuatro pasajes.
  \item \sevmedia\ \fileref{00\_resumen.tex} — Oración inicial muy cargada (tres subordinadas anidadas más una aposición). \textit{Recomendación:} partir la aposición de ``análisis estructural'' en oración aparte.
  \item \sevmedia\ Transversal — El cuarteto ``transparente, interactiva, verificable, coherente con el flujo profesional'' reaparece literal cuatro veces. \textit{Recomendación:} mantenerlo en hipótesis y conclusiones (donde el paralelismo es virtud) y parafrasear en el planteamiento y el campo de acción.
  \item \sevmedia\ \fileref{01\_introduccion.tex} — ``compañero de verificación'' es una personificación poco académica. \textit{(resuelto: sustituido por ``recurso de verificación'').}
  \item \sevmedia\ \fileref{01\_introduccion.tex} — Oración larga en la justificación del dominio 2D. \textit{(resuelto: dividida en dos).}
  \item \sevbaja\ Arranques de sección repetidos con ``abre / abre con''. \textit{Recomendación:} alternar con ``comienza con'', ``parte de''.
  \item \sevbaja\ \fileref{04\_resultados.tex} — ``concluyente'' + ``la evidencia más fuerte'' en oraciones contiguas, algo enfático. \textit{Recomendación:} atenuar.
  \item \sevbaja\ Transversal — Conviven guiones dobles de teclado y guiones largos tipográficos. \textit{Recomendación:} normalizar los incisos a guion largo.
  \item \sevbaja\ \fileref{04\_resultados.tex} (nota de Cook) — ``sobrepasada'' como sustantivo es informal. \textit{Recomendación:} ``ese exceso''.
\end{itemize}

\textbf{Veredicto.} Redacción sólida, de registro consistente y terminología impecable. La debilidad
real y accionable es la fatiga léxica por sobreuso de un puñado de fórmulas, sobre todo ``canal de
cálculo'' y el binomio ``caja negra\,/\,paso a paso''.

\subsection{Bibliografía y citación Vancouver}

\textbf{Fortalezas.} Integridad referencial perfecta: las 21 entradas se citan al menos una vez y todas
las invocaciones de \texttt{\textbackslash autocite} resuelven (cero claves huérfanas, cero entradas sin citar).
Estilo Vancouver correctamente configurado (biber, \texttt{numeric-comp}, \texttt{sorting=none}). Las
afirmaciones teóricas fuertes están ancladas donde corresponde y los datos propios del software no se
citan a fuentes externas.

\textbf{Hallazgos.}
\begin{itemize}
  \item \sevmedia\ \fileref{02\_marco\_teorico.tex} — Afirmaciones comparativas sobre productos de terceros sin fuente: VisualFEA (``cerrado, en inglés y de pago''), FEniCS y CALFEM en la tabla, ninguno en la bibliografía. \textit{Recomendación:} añadir una referencia mínima (manual/sitio o paper canónico) y citarla, o atenuar a una caracterización genérica sin nombres propios. No inventar DOIs.
  \item \sevmedia\ Brecha de literatura pedagógica: la base sobre enseñanza descansa en cuatro artículos del mismo journal, sin teoría educativa general. \textit{Recomendación:} incorporar una o dos referencias de fundamento (visualización/simulación interactiva en educación en ingeniería) para no sostener la justificación didáctica en una sola fuente.
  \item \sevbaja\ \fileref{02\_marco\_teorico.tex} — ``ED-Beams y ED-Frames'' se citan con el artículo de ED-Elas2D. \textit{Recomendación:} verificar que la fuente los nombre o reformular a ``programas hermanos del mismo grupo''.
  \item \sevbaja\ \fileref{00\_resumen.tex} — Las tasas teóricas de convergencia se respaldan con las fuentes del método MMS en vez de las de la teoría de aproximación. \textit{Recomendación (opcional):} añadir \fileref{strang2008analysis} a esa cita.
  \item \sevbaja\ \fileref{04\_resultados.tex} — El valor convergido de Cook ($\approx 23{,}95$) se atribuye a Cook (1974/2002); proviene de la literatura de \textit{benchmarking} posterior. \textit{Recomendación:} verificar o matizar ``ampliamente reportado en la literatura''.
\end{itemize}

\textbf{Veredicto.} Estado sólido: integridad referencial perfecta y Vancouver bien implementado. La
prioridad es cerrar las afirmaciones comparativas sin fuente sobre productos de terceros y ampliar
mínimamente la base de literatura pedagógica (dos a cuatro referencias, sin inventar datos).

\subsection{Higiene \LaTeX, figuras y tablas}

\textbf{Fortalezas.} Verificado de forma programática: 133 etiquetas, cero duplicados, cero referencias
cruzadas rotas; las 31 inclusiones de imagen resuelven contra la carpeta de figuras (ninguna falta);
tablas bien formadas (las columnas declaradas coinciden con los datos) y entornos matemáticos
balanceados. La nota de ``elaboración propia'' está presente y cubre el único activo de terceros.

\textbf{Hallazgos.}
\begin{itemize}
  \item \sevmedia\ \fileref{01\_introduccion.tex} — El Anexo F (memoria de cálculo paso a paso) no se referenciaba desde el cuerpo. \textit{(resuelto: añadido a la enumeración de anexos de la estructura).}
  \item \sevmedia\ \fileref{02\_marco\_teorico.tex} — La figura de elementos isoparamétricos (\fileref{fig:isoparametricos}) tiene leyenda y etiqueta pero no se cita en el texto. \textit{Recomendación:} citarla al introducir los elementos Q4/Q9.
  \item \sevbaja\ \fileref{07\_anexo\_memoria.tex} — La figura del modelo discretizado (\fileref{fig:mem-modelo}) no se cita. \textit{Recomendación:} citarla al abrir la sección de datos.
  \item \sevbaja\ \fileref{07\_anexo\_memoria.tex} — Cuatro tablas en \texttt{minipage} tienen etiqueta pero no se referencian. \textit{Recomendación:} citarlas en la prosa que las precede o aceptarlas como tablas de datos en línea.
  \item \sevbaja\ \fileref{07\_anexo\_memoria.tex} — Cuatro cajas \textit{overfull} ($\approx 53$~pt) por pares de \texttt{minipage} y dos matrices anchas. \textit{Recomendación:} bajar los anchos de \texttt{minipage} o envolver las matrices con \texttt{\textbackslash resizebox}.
  \item \sevbaja\ \fileref{preambulo.tex} — \texttt{\textbackslash headheight} de 12~pt dispara una advertencia repetida. \textit{Recomendación:} fijar \texttt{\textbackslash headheight} en 14~pt.
\end{itemize}

\textbf{Veredicto.} Higiene \LaTeX\ en muy buen estado: cero referencias rotas, cero figuras faltantes,
cero etiquetas duplicadas. La prioridad es integrar las figuras y el anexo huérfanos; el resto son
ajustes cosméticos.

\subsection{Anexos, reproducibilidad y aporte}

\textbf{Fortalezas.} Reproducibilidad sólida y verificable: el anexo de V\&V nombra cada guion con su
comando exacto y todos existen en el repositorio, junto con los datos crudos versionados y el
repositorio público citado. La memoria de cálculo (Anexo F) es numéricamente exacta y coherente con el
marco teórico (el equilibrio cierra, $\sigma_{\mathrm{VM}}$ es consistente). Los anexos no duplican el
cuerpo: aportan errores relativos, componentes adicionales y formatos de datos con detalle de producto.
El aporte queda demostrado.

\textbf{Hallazgos.}
\begin{itemize}
  \item \sevbaja\ \fileref{06\_anexos.tex} — La figura de la deformada de Cook usa una ruta con prefijo \texttt{figuras/} redundante (ya lo antepone \texttt{\textbackslash graphicspath}). \textit{Recomendación:} quitar el prefijo para alinearla con el resto.
  \item \sevbaja\ \fileref{06\_anexos.tex} — La deflexión del cuerpo (centroide) y la del punto A del anexo (a media altura) pueden leerse como contradicción. \textit{Recomendación:} aclarar que el punto A no está en el eje neutro.
  \item \sevbaja\ \fileref{06\_anexos.tex} (Anexo C) — Los listados de código no incluyen el solucionador (partición y resolución del sistema reducido). \textit{Recomendación (opcional):} añadir un fragmento corto o aclarar la omisión.
  \item \sevbaja\ \fileref{07\_anexo\_memoria.tex} — La configuración de apoyos del ejemplo canónico es atípica para un caso introductorio. \textit{Recomendación (opcional):} una frase aclarando que es el caso histórico de validación del proyecto.
\end{itemize}

\textbf{Veredicto.} Dimensión en muy buen estado: anexos completos, reproducibilidad excelente y memoria
de cálculo exacta que demuestra el aporte diferenciador. La única acción con prioridad clara es trivial
(normalizar una ruta de figura); el resto son refinamientos opcionales.

\section{Recomendaciones priorizadas}

La \autoref{tab:reco} consolida los hallazgos accionables en tres niveles de prioridad. Las
correcciones recaían en su mayoría sobre contenido preexistente; este informe las entrega para
decisión del autor, no las aplica (salvo las marcadas como ya incorporadas).

\begin{longtable}{@{}p{0.5cm}p{12cm}p{2.6cm}@{}}
\caption{Recomendaciones consolidadas por prioridad.}\label{tab:reco}\\
\toprule
\textbf{\#} & \textbf{Acción} & \textbf{Dimensión} \\
\midrule
\endfirsthead
\toprule
\textbf{\#} & \textbf{Acción} & \textbf{Dimensión} \\
\midrule
\endhead
\bottomrule
\endfoot
\multicolumn{3}{@{}l}{\textbf{\textcolor{cAlta}{Antes de la defensa (verificar y corregir)}}}\\[2pt]
1 & Corregir el error frente a Euler-Bernoulli: $1{,}82\%\to 1{,}85\%$ (o ajustar la referencia); verificar con \fileref{vv\_timoshenko}. & Rigor FEM \\
\midrule
\multicolumn{3}{@{}l}{\textbf{\textcolor{cMedia}{Mejoras de solidez (recomendadas)}}}\\[2pt]
2 & Podar la triple aparición de la lista de limitaciones; conservar en Resultados solo las que emergen de la evidencia medida. & Estructura \\
3 & Reencuadrar ``Resultados del software'' como cumplimiento del objetivo de cobertura y remitir a \fileref{tab:modulos}, sin re-describir módulos. & Estructura \\
4 & Diversificar el léxico sobre-usado: ``canal de cálculo'' ($\sim$28), ``caja negra'' (7) y ``paso a paso'' ($\sim$17). & Redacción \\
5 & Añadir una frase bisagra que distinga las variables numéricas (V\&V) de las pedagógicas (hipótesis de diseño). & Metodología \\
6 & Justificar el encuadre ``población/muestra'' sobre problemas-test y conectarlo con la validación pedagógica futura. & Metodología \\
7 & Dotar a PI-2 de un criterio de cumplimiento auditable (checklist etapa\,$\leftrightarrow$\,módulo). & Metodología \\
8 & Citar o atenuar las afirmaciones sobre VisualFEA, FEniCS y CALFEM; ampliar la base de literatura pedagógica (1--2 refs). & Bibliografía \\
9 & Citar la figura de elementos isoparamétricos, hoy huérfana. & \LaTeX/figuras \\
\midrule
\multicolumn{3}{@{}l}{\textbf{\textcolor{cBaja}{Pulido (opcional)}}}\\[2pt]
10 & Cifras de redondeo ($21{,}33\to 21{,}32$~GPa; nota de $0{,}2632\%$). & Rigor FEM \\
11 & ``(no evaluado; malla Q9)'' en \fileref{tab:resumen-q4q9}; sintetizar la doble validación de hipótesis. & Estructura \\
12 & Normalizar guiones largos; atenuar ``concluyente''; ``ese exceso'' por ``sobrepasada''. & Redacción \\
13 & Citar figuras y tablas huérfanas del Anexo F; corregir \textit{overfull}; fijar \texttt{\textbackslash headheight}. & \LaTeX/figuras \\
14 & Normalizar la ruta de la figura de Cook; aclarar punto A vs.\ centroide; (opcional) fragmento del solucionador en el Anexo C. & Anexos \\
\end{longtable}

\subsection*{Ya incorporado en esta intervención}

\begin{itemize}[nosep]
  \item Aclaración de que la tensión de von Mises se calcula en su forma plana y de que en deformación plana puede subestimar la tensión equivalente real (\fileref{02\_marco\_teorico.tex}).
  \item ``compañero de verificación'' $\to$ ``recurso de verificación'' (\fileref{01\_introduccion.tex}).
  \item División de la oración larga en la justificación del dominio 2D (\fileref{01\_introduccion.tex}).
  \item Referencia al Anexo F (memoria de cálculo) en la enumeración de anexos (\fileref{01\_introduccion.tex}).
  \item Reanclaje del problema general tras la corrección de ED-Elas2D, y citas localizadas en cada bloqueo (revisión de la intervención).
\end{itemize}

\subsection*{Decisiones de autor (legítimas y defendibles)}

Las siguientes observaciones no son errores: corresponden a decisiones deliberadas que conviene saber
defender ante el tribunal, no necesariamente cambiar.

\begin{itemize}[nosep]
  \item \textbf{Tres capítulos con el diseño metodológico integrado} en el capítulo de diseño: defendible como unidad propia de un trabajo tecnológico; basta una frase justificativa.
  \item \textbf{``Población/muestra'' sobre problemas-test} y no sobre sujetos: coherente con una tesis tecnológica cuyo objeto evaluable es el artefacto; la población humana corresponde a la validación pedagógica futura.
  \item \textbf{Reporte de $0{,}2633\%$} (precisión completa del motor) frente al $0{,}2632\%$ de los valores redondeados de tabla: admisible.
  \item \textbf{Configuración de apoyos del ejemplo canónico}: es el caso histórico de validación del proyecto, internamente consistente.
\end{itemize}

\section{Conclusión del auditor}

La tesis está en un estado maduro y apto para la defensa. Su columna vertebral —la cadena
problema\,--\,objetivo\,--\,hipótesis\,--\,verificación\,--\,conclusiones— es coherente y se sostiene con
evidencia numérica trazable y un marco metodológico bien preparado para el tribunal. \textbf{La
auditoría no halló ningún defecto que comprometa la validez de los resultados:} de los 43 hallazgos,
los dos de severidad alta son de estilo y el único error aritmético genuino es de fácil corrección.

El grueso del trabajo restante es editorial: podar la redundancia transversal, diversificar un puñado
de expresiones sobre-usadas y reforzar dos flancos metodológicos con frases puntuales. Atendidas la
corrección aritmética y las mejoras de solidez de prioridad media, el documento quedará en condiciones
óptimas. Este informe funciona como hoja de ruta; la decisión sobre cada punto corresponde al autor.

\end{document}
